\documentclass[twocolumn]{aastex631}

\usepackage{amsmath}
\usepackage{multirow}
\usepackage{natbib}
\usepackage{graphicx} 
\usepackage{aas_macros}

\begin{document}

Higher-order scalar-tensor theories, such as Degenerate Higher-Order Scalar-Tensor (DHOST) theories, typically suffer from Ostrogradsky instabilities due to the presence of higher-order derivatives. While the Type Ia DHOST subclass is classically designed to circumvent this problem by satisfying a specific degeneracy condition that eliminates the unwanted ghost degree of freedom, the fundamental origin of this stability has remained underexplored. This paper proposes and demonstrates that this classical degeneracy is not an ad-hoc fine-tuning but rather the direct manifestation of a novel, non-linear, field-dependent gauge symmetry inherent to the Type Ia DHOST Lagrangian. Our investigation first establishes the classical foundation by explicitly deriving the general DHOST degeneracy condition and confirming its satisfaction within Type Ia theories, thereby ensuring the absence of ghosts and positive propagation speeds for scalar and tensor modes. We then crucially demonstrate that this degeneracy condition is mathematically equivalent to the requirement for invariance under a specific field-dependent gauge transformation, for which we explicitly compute the transformation functions. Leveraging this discovery, we employ the path integral formalism to derive the associated Ward-Takahashi identities, which are subsequently used to analyze the quantum stability of the theory. We show that these identities impose stringent algebraic constraints on the structure of one-loop radiative counterterms, thereby preventing the generation of problematic higher-derivative ghost terms and ensuring the theory's robustness against radiative instabilities. Finally, by examining the behavior of this degeneracy-induced symmetry in the Beyond Horndeski and Horndeski limits, we reveal that while it persists and provides quantum protection in the former, it becomes trivial in the latter, clarifying the distinct stability mechanisms across the hierarchy of scalar-tensor theories.

\end{document}
                